\documentclass{beamer}
\usepackage[latin1]{inputenc}
\beamertemplateshadingbackground{red!10}{blue!10}
%\usepackage{fancybox}
\usepackage{epsfig}
\usepackage{verbatim}
\usepackage{url}
%\usepackage{graphics}
%\usepackage{xcolor}
\usepackage{fancybox}
\usepackage{moreverb}
%\usepackage[all]{xy}
\usepackage{listings}
\usepackage{filecontents}
\usepackage{graphicx}

\lstset{
  language=Lisp,
  basicstyle=\scriptsize\ttfamily,
  keywordstyle={},
  commentstyle={},
  stringstyle={}}

\def\inputfig#1{\input #1}
\def\inputeps#1{\includegraphics{#1}}
\def\inputtex#1{\input #1}

\inputtex{logos.tex}

%\definecolor{ORANGE}{named}{Orange}

\definecolor{GREEN}{rgb}{0,0.8,0}
\definecolor{YELLOW}{rgb}{1,1,0}
\definecolor{ORANGE}{rgb}{1,0.647,0}
\definecolor{PURPLE}{rgb}{0.627,0.126,0.941}
\definecolor{PURPLE}{named}{purple}
\definecolor{PINK}{rgb}{1,0.412,0.706}
\definecolor{WHEAT}{rgb}{1,0.8,0.6}
\definecolor{BLUE}{rgb}{0,0,1}
\definecolor{GRAY}{named}{gray}
\definecolor{CYAN}{named}{cyan}

\newcommand{\orchid}[1]{\textcolor{Orchid}{#1}}
\newcommand{\defun}[1]{\orchid{#1}}

\newcommand{\BROWN}[1]{\textcolor{BROWN}{#1}}
\newcommand{\RED}[1]{\textcolor{red}{#1}}
\newcommand{\YELLOW}[1]{\textcolor{YELLOW}{#1}}
\newcommand{\PINK}[1]{\textcolor{PINK}{#1}}
\newcommand{\WHEAT}[1]{\textcolor{wheat}{#1}}
\newcommand{\GREEN}[1]{\textcolor{GREEN}{#1}}
\newcommand{\PURPLE}[1]{\textcolor{PURPLE}{#1}}
\newcommand{\BLACK}[1]{\textcolor{black}{#1}}
\newcommand{\WHITE}[1]{\textcolor{WHITE}{#1}}
\newcommand{\MAGENTA}[1]{\textcolor{MAGENTA}{#1}}
\newcommand{\ORANGE}[1]{\textcolor{ORANGE}{#1}}
\newcommand{\BLUE}[1]{\textcolor{BLUE}{#1}}
\newcommand{\GRAY}[1]{\textcolor{gray}{#1}}
\newcommand{\CYAN}[1]{\textcolor{cyan }{#1}}

\newcommand{\reference}[2]{\textcolor{PINK}{[#1~#2]}}
%\newcommand{\vect}[1]{\stackrel{\rightarrow}{#1}}

% Use some nice templates
\beamertemplatetransparentcovereddynamic

\newcommand{\A}{{\mathbb A}}
\newcommand{\degr}{\mathrm{deg}}

\title{Fast, Maintainable, and Portable Sequence Functions}

\author{Ir�ne Durand \& Robert Strandh}
\institute{
LaBRI, University of Bordeaux
}
\date{April, 2017}

%\inputtex{macros.tex}


\begin{document}
\frame{
\resizebox{3cm}{!}{\includegraphics{Logobx.pdf}}
\hfill
\resizebox{1.5cm}{!}{\includegraphics{labri-logo.pdf}}
\titlepage
\vfill
\small{European Lisp Symposium, Brussels, Belgium \hfill ELS2017}
}

\setbeamertemplate{footline}{
\vspace{-1em}
\hspace*{1ex}{~} \GRAY{\insertframenumber/\inserttotalframenumber}
}

\frame{
\frametitle{Context: The \sicl{} project}

https://github.com/robert-strandh/SICL

Several objectives:

\begin{itemize}
\item Create high-quality \emph{modules} for implementors of
  \commonlisp{} systems.
\item Improve existing techniques with respect to algorithms and data
  structures where possible.
\item Improve readability and maintainability of code.
\item Improve documentation.
\item Ultimately, create a new implementation based on these modules.
\end{itemize}
}


\frame{
\frametitle{The challenge}

  \begin{itemize}
  \item Good performance requires specialization on element type (for
    vector), test functions, key functions, and values of certain
    parameters such as \texttt{:end} and \texttt{from-end} (for
    lists).
  \item A separate function for each case makes the source code huge
    and hard to test.
  \item Using macros to generate separate functions makes the code
    hard to understand.
  \end{itemize}
}

\frame{
\frametitle{Previous work}

At ELS 2015 in London, we presented a technique for processing lists
from the end (when \texttt{:from-end} is true).  The present work uses
that technique to obtain one special case.
}

\frame{
\frametitle{Previous work}

\begin{itemize}
\item \ecl{} Uses macros similar to ours, but does not specialize
\item \sbcl{} Uses a mixed approach.  It Requires the caller to specify
  type in order to be fast.
\item The CLISP sequence functions are implemented in C.
\end{itemize}

}

\frame{
\frametitle{Our technique}

\begin{itemize}
\item We use macros for each required specialization.
\item The macros duplicate the body so that the compiler can
  specialize each case.
\end{itemize}
}

\frame[containsverbatim]{
\frametitle{Example: handling \texttt{:end}}

To illustrate our technique, we use a simplified version of
\texttt{find}.

\begin{verbatim}
(defun find (item list end)
  (loop for element in list
        for index from 0
        until (and end (>= index end))
        when (eql item element)
          return element))
\end{verbatim}

Whether \texttt{end} is \texttt{nil} or not is loop invariant, but
obviously not known at compile time, so a typical compiler can not
do anything with this information.
}
\frame[containsverbatim]{
\frametitle{Example: handling \texttt{:end}}

To get better performance, we could do this:

\begin{verbatim}
(defun find (item list end)
  (if end
      (loop for element in list
            for index from 0
            until (and end (>= index end))
            when (eql item element)
              return element)
      (loop for element in list
            when (eql item element)
              return element)))
\end{verbatim}

But now we have code duplication which makes this version less
maintainable and harder to test.
}

\frame[containsverbatim]{
\frametitle{Example: handling \texttt{:end}}

Our technique wraps the loop in a macro call like this:

\begin{verbatim}
(defun find (item list end)
  (with-end
    (loop for element in list
          for index from 0
          until (and end (>= index end))
          when (eql item element)
            return element)))
\end{verbatim}

}

\frame[containsverbatim]{
\frametitle{Example: handling \texttt{:end}}

The macro is defined like this:

\begin{verbatim}
(defmacro with-end (&body body)
  `(if end
       (progn ,@body)
       (progn ,@body)))
\end{verbatim}

}

\frame{
\frametitle{Future work}

\begin{itemize}
\item Determine policy.
\end{itemize}

}

\frame{
  \frametitle{Acknowledgments}

We would like to thank Bart Botta, Pascal Bourguignon, and Philipp
Marek for providing valuable feedback on early versions of this paper.
}

\frame{
\frametitle{Thank you}

Questions?
}

%% \frame{\tableofcontents}
%% \bibliography{references}
%% \bibliographystyle{alpha}

\end{document}
